\documentclass[a4paper,10.5pt]{article}
\usepackage[utf8]{inputenc}
\usepackage[T1]{fontenc}
\usepackage[french]{babel}
\usepackage[left=1cm,right=1cm,top=.5cm,bottom=0cm]{geometry}
\usepackage{enumitem}
\usepackage{hyperref}
\usepackage{xcolor}
\usepackage{titlesec}
\usepackage{fontawesome5}
\usepackage{lmodern}
\usepackage[normalem]{ulem}

% --- Couleurs Sopra Steria (Rouge corporate) ---
\definecolor{primary}{RGB}{190, 30, 45}
\definecolor{secondary}{RGB}{80, 80, 80}

% --- Config Liens ---
\hypersetup{colorlinks=true, linkcolor=primary, filecolor=primary, urlcolor=primary}

% --- Titres ---
\titleformat{\section}{\large\bfseries\color{primary}\uppercase}{}{0em}{}[\titlerule]
\titlespacing{\section}{0pt}{8pt}{5pt}

\begin{document}

% --- EN-TÊTE ---
\begin{center}
    {\Large \textbf{Loïc Aron MBASSI EWOLO}} \\ \vspace{4pt}
    \textbf{\large Alternance Développeur Full Stack IA} \\
    \textit{\small Disponible dès septembre 2026} \\ \vspace{4pt}
    \small
    \faMapMarker\ Cergy, France (Mobile nationale) \quad
    \faPhone\ (+33) 7 59 52 33 98 \quad
    \faEnvelope\ \href{mailto:loicaron.mbassiewolo@ensea.fr}{loicaron.mbassiewolo@ensea.fr} \\ \vspace{2pt}
    \faLinkedin\ \href{https://www.linkedin.com/in/loic-mbassi/}{linkedin.com/in/loic-mbassi} \quad
    \faGithub\ \href{https://github.com/Nameless0l}{github.com/Nameless0l} \quad
    \faGlobe\ \href{https://mbassiloic.tech}{mbassiloic.tech}
\end{center}

\vspace{-6pt}

% --- PROFIL ---
\noindent\small{Élève-ingénieur double diplôme ENSEA/Polytechnique Yaoundé avec une expérience concrète en \textbf{développement d'agents IA (RAG, LangChain, FastAPI)} et en \textbf{NLP/Transformers (BERT, GPT-2, HuggingFace)}. Stage de recherche au MILA (Institut Québécois d'IA) sur l'interprétabilité des modèles. +3 ans d'expérience en développement Full Stack (Python, PHP, JS). Habitué au travail en Agile, Docker et Git.}

% --- FORMATION ---
\section{Formation}
\begin{itemize}[leftmargin=*, label=\tiny$\bullet$, nosep]
    \item \textbf{ENSEA -- École Nationale Supérieure de l'Électronique et de ses Applications} \hfill \textit{2025 -- 2027} \\
    \small{Double diplôme d'Ingénieur -- Systèmes Embarqués, Signal, IA. Cursus : Deep Learning, traitement du signal, architecture logicielle.}
    \item \textbf{École Polytechnique de Yaoundé (1\textsuperscript{ère} école d'ingénieur CEMAC)} \hfill \textit{2021 -- 2025} \\
    \small{Diplôme d'Ingénieur de Conception (Master 2) : \textbf{Mathématiques Appliquées}, Génie Logiciel, Algorithmique.}
\end{itemize}

% --- COMPÉTENCES ---
\section{Compétences Techniques}
\begin{itemize}[leftmargin=*, nosep]
    \item \small{\textbf{IA / NLP :} PyTorch, TensorFlow/Keras, LangChain, HuggingFace (Transformers, Speech), RAG, BERT, GPT-2, OpenCV.}
    \item \small{\textbf{Backend \& APIs :} Python (Expert), FastAPI, PHP/Laravel, Node.js/NestJS, REST APIs, OAuth2/JWT.}
    \item \small{\textbf{Bases de données :} PostgreSQL, MongoDB, MySQL, Redis, bases vectorielles.}
    \item \small{\textbf{DevOps \& Outils :} Git/GitLab, Docker, CI/CD, Linux (Ubuntu), Jira, tests unitaires, documentation technique.}
    \item \small{\textbf{Méthodologies :} Agile/Scrum, Clean Code, code review. \textbf{Langues :} Français (natif), Anglais (professionnel/technique).}
\end{itemize}

% --- PROJETS IA ---
\section{Projets IA \& Agents}
\begin{itemize}[leftmargin=*, label={-}]

\item \textbf{Système Multi-Agents Agricole (RAG + LLM)} \hfill \textit{2024 -- 2025} \\
\textit{Projet de Recherche} \hfill \textit{Python, Google ADK, LangChain, RAG, FastAPI, Docker}
\vspace{-2pt}
    \begin{itemize}[leftmargin=0.4cm, nosep, label=\tiny$\bullet$]
        \item \small{Architecture de 4 agents IA collaboratifs utilisant Gemini (LLM) et \textbf{RAG} pour l'enrichissement contextuel.}
        \item \small{Orchestration via Google ADK avec résolution de conflits entre agents et fusion de recommandations.}
        \item \small{Intégration d'APIs externes (OpenWeather, marchés) via \textbf{FastAPI}. Déploiement Docker reproductible.}
    \end{itemize}
\vspace{3pt}

\item \textbf{YembaAudio -- NLP sur Langue Peu Dotée} \hfill \textit{2025} \\
\textit{Projet Académique -- Préservation linguistique} \hfill \textit{Python, HuggingFace, TensorFlow, Next.js}
\vspace{-2pt}
    \begin{itemize}[leftmargin=0.4cm, nosep, label=\tiny$\bullet$]
        \item \small{App web d'apprentissage d'une langue locale à partir d'un \textbf{modèle pré-entraîné HuggingFace}.}
        \item \small{Reconnaissance vocale (Speech Recognition) sur langue peu dotée avec corpus custom.}
    \end{itemize}
\vspace{3pt}

\item \textbf{Women Safe -- Détection de Contenus Sexistes (NLP)} \hfill \textit{2023} \\
\textit{Compétition MLPC} \hfill \textit{Python, BERT, GPT-2, Transformers}
\vspace{-2pt}
    \begin{itemize}[leftmargin=0.4cm, nosep, label=\tiny$\bullet$]
        \item \small{Modèle NLP de détection de misogynie en ligne basé sur \textbf{BERT et GPT-2} (fine-tuning Transformers).}
        \item \small{Étude des biais éthiques et analyse des performances sur des corpus multilingues.}
    \end{itemize}
\vspace{3pt}

\end{itemize}

% --- EXPÉRIENCES RECHERCHE & DEV ---
\section{Recherche \& Développement Full Stack}
\begin{itemize}[leftmargin=*, label={-}]

\item \textbf{Assistant de Recherche @ MILA (Institut Québécois d'IA)} \hfill \textit{Juin 2025 -- Sept 2025} \\
\textit{Remote -- Interprétabilité des Modèles} \hfill \textit{Python, PyTorch, TensorFlow}
\vspace{-2pt}
    \begin{itemize}[leftmargin=0.4cm, nosep, label=\tiny$\bullet$]
        \item \small{Étude du phénomène de Grokking (généralisation tardive) sur des MLP avec PyTorch.}
        \item \small{Reproduction d'expériences SOTA, analyse mathématique de la convergence et rapports de synthèse.}
    \end{itemize}
\vspace{3pt}

\item \textbf{Développeur Web Full Stack @ SOSDOCTA} \hfill \textit{Fév 2022 -- Mars 2024} \\
\textit{Remote (Belgique/Canada) -- Télémédecine} \hfill \textit{Laravel, PHP, MySQL, JS, OAuth2/JWT}
\vspace{-2pt}
    \begin{itemize}[leftmargin=0.4cm, nosep, label=\tiny$\bullet$]
        \item \small{Développement complet d'une plateforme de télémédecine : backend REST, frontend, paiement sécurisé.}
        \item \small{Maintenance évolutive 2 ans en autonomie. Gestion données sensibles (RGPD).}
    \end{itemize}
\vspace{3pt}

\item \textbf{Laravel API Generator -- Outil Open Source} \hfill \textit{2024 -- Présent} \\
\textit{Projet Personnel} \hfill \textit{PHP, Python, XML, LLM Integration, Packagist}
\vspace{-2pt}
    \begin{itemize}[leftmargin=0.4cm, nosep, label=\tiny$\bullet$]
        \item \small{Outil de génération automatique d'architecture backend depuis UML/XML. Intégration LLM pour assistance.}
        \item \small{Open source : +30 étoiles GitHub, déployé sur Packagist. Documentation technique complète.}
    \end{itemize}
\vspace{3pt}

\end{itemize}

% --- DISTINCTIONS ---
\section{Leadership \& Distinctions}
\begin{itemize}[leftmargin=*, label=\tiny$\bullet$, nosep]
    \item \small{\textbf{1\textsuperscript{er} Prix} IndabaX Africa 2025 (IA Santé : déploiement API Flask, dashboard Streamlit).}
    \item \small{\textbf{Lead AI Cell} Polytechnique Yaoundé : workshops (Python, ML, fine-tuning LLM), collab. Google DeepMind.}
    \item \small{\textbf{Certifications :} Supervised ML (DeepLearning.AI/Stanford), Python for Data Science (IBM/Coursera).}
\end{itemize}

\end{document}
